\section{Introduction}

\ac{CCA} are a cornerstone of modern network technology, serving as essential mechanisms to regulate data flow and optimize the utilization of network resources. 
\ac{CCA} is an key topic to the internet as a whole and its being researched since over 40 years so far.
Althought this is a long time, there are still improvements done and research is also actively done.
The importance of \ac{CCA} is that it can be the driving factor to achieve low latency and high thoughtput as well as helping to improve other characteristics of the internet.

This paper provides a comprehensive overview of contemporary \acs{CCA}, examining their operational logic and comparing them through the lens of their design objectives.
To establish a foundational understanding, Section \ref{sec:Background} presents the necessary theoretical background and common performance metrics. 
This section explores established protocols such as TCP Cubic, TCP Vegas, and FAST TCP, while detailing the metrics used to evaluate their efficacy.

In Section \ref{sec:Approach}, we survey recent advancements in the field, specifically focusing on Bottleneck Bandwidth and Round-trip propagation time (BBR) and Copa, alongside both of their extensions. 
We also highlight several other emerging algorithms that represent the current state of the art.
Finally, Section \ref{sec:Evaluation} provides a comparative analysis of these algorithms, evaluating them based on the performance data and architectural goals provided by their respective authors. 
We conclude by situating our findings within the broader academic context, referencing existing literature that offers deeper empirical evaluations of these protocols.
We conclude the paper in Section \ref{sec:Summary} by summarizing our findings and synthesizing the key takeaways from our comparison. 
Furthermore, we provide an outlook in Section \ref{sec:Outlook} on the future of the field, identifying potential trends and remaining challenges in the evolution of congestion control.
